\begin{center}
\begin{tikzpicture}[scale=1.25, line cap=round, line join=round]
  \definecolor{trigblue}{RGB}{30,110,190}
  \coordinate (A) at (0,0);
  \coordinate (B) at (3.4,0);
  \coordinate (C) at (3.4,2.0);
  \draw[very thick] (A)--(B)--(C)--cycle;
  \draw[line width=2.2pt,trigblue] (A)--(C);
  \draw[line width=2.2pt,gray!65] (B)--(C);
  \draw (3.1,0)--(3.1,0.3)--(3.4,0.3);
  \draw[trigblue] (0.55,0) arc[start angle=0,end angle=30.5,radius=0.55];
  \node[trigblue] at (0.78,0.22) {$\theta$};
  \node[above left,trigblue] at ($(A)!0.5!(C)$) {hypotenuse};
  \node[right] at ($(B)!0.5!(C)$) {opposite};
  \node[below] at (1.7,-0.55) {$\csc\theta=\text{hypotenuse}/\text{opposite}$};
\end{tikzpicture}
\end{center}


\begin{definitionbox}{Definition}
\begin{definition}[Right-triangle cosecant]
\label{def:right-triangle-cosecant}
Let $T$ be a right triangle and let $\theta$ be one of its acute angles. Define
\[
\csc\theta=\frac{\operatorname{hyp}(T)}{\operatorname{opp}(T,\theta)}.
\]
\end{definition}
\end{definitionbox}
\begin{remark*}[Standard quantified statement]
For every admissible input satisfying the hypotheses of the statement, the
displayed definition or assertion holds in the stated domain.
\end{remark*}

\begin{remark*}[Interpretation]
The statement records the named trigonometric object or identity together with
its intended domain of use.
\end{remark*}

\begin{dependencies}
\begin{itemize}
  \item \hyperref[def:euclid-triangle-angle-classes]{Right-angled triangle}
  \item \hyperref[def:euclid-right-angle-and-perpendicular]{Right angle and perpendicular}
  \item \hyperref[def:euclid-acute-angle]{Acute angle}
\end{itemize}
\end{dependencies}

\input{volume-i/book-geometry/trigonometry/foundations/figure-extended-cosecant}

\begin{definitionbox}{Definition}
\begin{definition}[Extended cosecant]
\label{def:extended-cosecant}
Where $\sin\theta\neq 0$, define
\[
  \csc\theta=\frac{1}{\sin\theta}.
\]
\end{definition}
\end{definitionbox}
\begin{remark*}[Standard quantified statement]
For every admissible input satisfying the hypotheses of the statement, the
displayed definition or assertion holds in the stated domain.
\end{remark*}

\begin{remark*}[Interpretation]
The statement records the named trigonometric object or identity together with
its intended domain of use.
\end{remark*}

\begin{dependencies}
\begin{itemize}
  \item \hyperref[def:unit-circle-sine]{Unit-circle sine}
\end{itemize}
\end{dependencies}
